\documentclass[conference]{IEEEtran}
\usepackage{amsmath,amssymb}
\usepackage{booktabs}
\usepackage{hyperref}
\usepackage{bm}
\usepackage{multirow}

\title{PPO-Based Autonomous Drone Racing\\through Gate Sequences in Isaac Sim}
\author{
\IEEEauthorblockN{Jinyuan Zhang, Ningze Zhong}
\IEEEauthorblockA{\{jinyuanz, zhong666\}@seas.upenn.edu}
}

\begin{document}
\maketitle

\section{Overview}

We train a reinforcement learning policy using Proximal Policy Optimization (PPO)~\cite{schulman2017ppo} to fly a simulated Crazyflie quadrotor through a seven-gate racing circuit in NVIDIA Isaac Sim. The track contains a challenging powerloop segment in which two consecutive gates share the same yaw orientation ($\psi_{g_2}=\psi_{g_3}=+\frac{\pi}{2}$), high-altitude gates ($z=2.0$\,m), and a chicane. The layout was intended to induce a vertical loop, although the final policy often bypasses that maneuver. Our design draws on prior work in autonomous drone racing~\cite{song2021autonomous, kaufmann2023champion, song2023reaching} and reactive aerobatic flight~\cite{han2025reactive}. We also explored the Eureka framework~\cite{ma2024eureka} for reward tuning. The final system uses four elements: (i)~a compact reward with dense progress shaping, sparse gate-pass bonuses, and a per-step time penalty, (ii)~a 36-dimensional body-frame observation space, (iii)~a five-stage reset curriculum, and (iv)~domain randomization of dynamics parameters. The best policy, trained with 4096 parallel environments, completes three laps in 18.26\,s.

\section{Reward Function Design}

We initially adopted the projection-based progress reward from Song et al.~\cite{song2021autonomous}, which measures advancement along the line segment between adjacent gate centers. That formulation was sufficient for standard gate sequences, but it did not produce successful powerloop traversals. We then redesigned the reward by combining insights from prior work~\cite{kaufmann2023champion, song2023reaching, han2025reactive} and briefly explored Eureka~\cite{ma2024eureka} for reward tuning. The final reward has three active terms:
\begin{equation}\label{eq:reward}
  r(t) = w_{\text{prog}}\, r_{\text{prog}}(t) + w_{\text{gate}}\, r_{\text{gate}}(t) + w_{\text{time}}\, r_{\text{time}}.
\end{equation}
Upon terminal failure, the reward is overridden by a death cost $c_{\text{death}}$. The active reward weights in the current implementation are listed in Table~\ref{tab:reward}.

\begin{table}[h]
\centering
\caption{Active reward component weights.}
\label{tab:reward}
\begin{tabular}{@{}lrc@{}}
\toprule
\textbf{Component} & \textbf{Symbol} & \textbf{Weight} \\
\midrule
Gate-pass bonus (sparse) & $w_{\text{gate}}$ & $55.0$ \\
Progress reward (dense) & $w_{\text{prog}}$ & $15.0$ \\
Time penalty (per-step) & $w_{\text{time}}$ & $-0.08$ \\
Death cost (terminal) & $c_{\text{death}}$ & $-80.0$ \\
\bottomrule
\end{tabular}
\end{table}

\textbf{Global progress reward.}
The primary dense shaping signal measures the drone's advancement along the racing line. We define global progress as
\begin{equation}\label{eq:progress}
  G(t) = n_{\text{passed}}(t) + 1 - \tanh\!\left(\frac{d(t)}{\sigma_d}\right),
\end{equation}
where $n_{\text{passed}}(t) \in \mathbb{Z}_{\geq 0}$ is the cumulative number of gates passed, $d(t) = \|\mathbf{p}_{\text{drone}}(t) - \mathbf{p}_{g_i}\|_2$ is the Euclidean distance to the current target gate center $\mathbf{p}_{g_i}$, and $\sigma_d = 3.0$\,m is a distance scale. The per-step increment is
\begin{equation}\label{eq:rprog}
  r_{\text{prog}}(t) = \text{clamp}\!\bigl(G(t) - G(t{-}1),\; -1,\; 1\bigr).
\end{equation}
The $\tanh$ normalization keeps the signal smooth and bounded. On steps where a gate is crossed, $r_{\text{prog}}$ is set to zero to avoid double-counting with the gate-pass bonus. For tight same-direction gate pairs in the powerloop segment, negative progress increments are clipped to zero.

\textbf{Gate-pass bonus.}
A sparse reward $r_{\text{gate}}(t) \in \{0, 1\}$ fires when the drone crosses a gate. In the gate-local frame $\mathcal{F}_g$ with origin at the gate center and the normal $\hat{\mathbf{x}}_g$ pointing opposite to the approach direction, crossing is detected as
\begin{equation}\label{eq:gate_cross}
  x_g(t{-}1) > 0.1 \;\;\wedge\;\; x_g(t) \leq 0.1,
\end{equation}
subject to the aperture constraint $|y_g(t)| < \tfrac{s}{2}$ and $|z_g(t)| < \tfrac{s}{2}$, where $s = 1.0$\,m is the gate side length. Upon a valid crossing, the waypoint index advances as $i \leftarrow (i + 1) \bmod N$, where $N=7$.

\textbf{Time penalty.}
A constant per-step cost $r_{\text{time}} = 1.0$ (scaled by $w_{\text{time}}=-0.08$) discourages hovering and implicitly incentivizes speed without an explicit velocity reward, which in early experiments caused the drone to fly fast in unproductive directions.

\textbf{Ablated reward terms.}
During development we explored four additional reward components, all disabled in the final configuration. A \emph{speed bonus} $r_{\text{speed}} = \tanh(\|\mathbf{v}\|/3.0)$ encouraged unproductive fast flight. An \emph{action smoothness penalty} $r_{\text{smooth}} = \|\mathbf{a}(t) - \mathbf{a}(t{-}1)\|_2^2$ reduced the aggressive motions needed for the powerloop. A \emph{crash penalty} triggered when the contact sensor measured a nonzero net force on the body after a 100-step warm-up. This made the policy too conservative near gate edges. An \emph{entry half-plane reward} gave a one-time bonus when the drone returned to the correct approach side ($x_g > 0.15$\,m) before a tight same-direction gate pair. Mid-loop resets were more effective, so this term was removed.

\section{Observation Space Design}

Following Kaufmann et al.~\cite{kaufmann2023champion} and Song et al.~\cite{song2023reaching}, the policy receives a 36-dimensional observation vector $\mathbf{o}(t) \in \mathbb{R}^{36}$:
\begin{equation}\label{eq:obs}
  \mathbf{o}(t) = \bigl[\,\mathbf{v}^b,\;\text{vec}(\mathbf{R}_{bw}),\;\mathbf{c}^b_{i,1{:}4},\;\mathbf{c}^b_{i{+}1,1{:}4}\,\bigr].
\end{equation}
The components are summarized in Table~\ref{tab:obs}.

\begin{table}[h]
\centering
\caption{Observation vector components ($\dim = 36$).}
\label{tab:obs}
\begin{tabular}{@{}llcc@{}}
\toprule
\textbf{Group} & \textbf{Component} & \textbf{Dim} & \textbf{Frame} \\
\midrule
\multirow{2}{*}{Ego-state}
  & Linear velocity $\mathbf{v}^b$ & 3 & Body \\
  & Rotation matrix $\text{vec}(\mathbf{R}_{bw})$ & 9 & B$\to$W \\
\midrule
\multirow{2}{*}{Gate info}
  & Current gate corners $\mathbf{c}^b_{i,1{:}4}$ & 12 & Body \\
  & Next gate corners $\mathbf{c}^b_{i{+}1,1{:}4}$ & 12 & Body \\
\midrule
  & \textbf{Total} & \textbf{36} & \\
\bottomrule
\end{tabular}
\end{table}

\textbf{Ego-state} (12D). The linear velocity in the body frame $\mathbf{v}^b \in \mathbb{R}^3$ is read from the robot's center-of-mass velocity. The body-to-world rotation matrix $\mathbf{R}_{bw} \in SO(3)$ is converted from the world-frame quaternion $\mathbf{q}_w = (q_w, q_x, q_y, q_z)$ and flattened row-major into $\text{vec}(\mathbf{R}_{bw}) \in \mathbb{R}^9$. We use the full rotation matrix rather than the quaternion to avoid sign ambiguity and to provide a smoother representation.

\textbf{Gate geometry} (24D). For each of the current ($i$) and next ($i{+}1 \bmod N$) target gates, the four corner positions in the body frame are computed as
\begin{equation}\label{eq:corners}
  \mathbf{c}^b_{j,k} = \mathbf{R}_{bw}^\top \bigl(\mathbf{c}^w_{j,k} - \mathbf{p}_{\text{drone}}\bigr), \quad k=1,\dots,4,
\end{equation}
where $\mathbf{c}^w_{j,k}$ are the world-frame corner positions and $\mathbf{p}_{\text{drone}}$ is the drone position. Each gate is a $1\,\text{m}\times 1\,\text{m}$ square with local corners at $\mathbf{c}_k^{\text{local}} \in \{(0, \pm 0.5, \pm 0.5)\}$. The corners implicitly encode gate position, orientation, distance, and size. Absolute world position is omitted to reduce overfitting to track-specific coordinates. The implementation also supports an optional privileged critic that appends body-frame angular velocity and world position, yielding a 42-dimensional critic observation during training.

\section{Reset Strategy and Curriculum Learning}

Since the progress reward alone could not produce useful powerloop states, early on-policy rollouts rarely reached the gate 2$\rightarrow$3 segment. We therefore used curriculum learning to reshape the initial-state distribution over training. Early stages focus on stable starts and the first gate. Later stages unlock harder mid-track states while keeping a nonzero fraction of gate-0 starts so the policy does not forget end-to-end racing from the beginning of the track. Inspired by Han et al.~\cite{han2025reactive}, we designed a five-stage curriculum indexed by the training progress ratio $\rho = k / K$, where $k$ is the current iteration and $K=1500$ for the reported run. The stage schedule is listed in Table~\ref{tab:curriculum}.

At each stage, the curriculum controls two quantities. First, $U$ denotes the number of unlocked reset targets, so resets may start from gates $\{0,\dots,U-1\}$. Second, $p_0$ denotes the probability of resetting at gate 0. We decrease $p_0$ gradually from 1.0 to 0.35 rather than driving it to zero, because later training still benefits from revisiting takeoff and early-gate approach behavior.

\textbf{Position randomization.}
At each gate $g_i$ with yaw $\psi_i$, the drone is placed 0.5--3.0\,m behind the gate along its normal, with lateral offset $\pm 1.0$\,m and height noise $\pm 0.3$\,m, then rotated into world coordinates via $\psi_i$. The orientation faces the gate with additive noise: $\Delta\phi, \Delta\theta \sim \mathcal{U}(-0.1, 0.1)$\,rad and $\Delta\psi \sim \mathcal{U}(-0.2, 0.2)$\,rad. Mid-track resets initialize forward velocity $v_0 \sim \mathcal{U}(1.5, 4.0)$\,m/s along the gate approach direction $-\hat{\mathbf{n}}_{g_i}$.

\textbf{Powerloop mid-loop resets.}
After $\rho > 0.10$, a separate apex-reset branch is activated for the powerloop region, i.e., the gate 2$\rightarrow$3 segment. Up to 30\% of all training resets are reassigned to this branch. A continuous phase $\alpha \sim \mathcal{U}(0, 1)$ parameterizes the sampled loop state, where $\alpha=0$ corresponds to loop-entry states just after gate 2 and $\alpha=1$ corresponds to near-apex states before descending toward gate 3. The sampled height is $z = 0.8 + 1.5\alpha \pm 0.15$\,m, the sampled loop pitch is $\vartheta_{\text{loop}} = \alpha \cdot \vartheta_{\max}$, where $\vartheta_{\max} = \frac{2\pi}{3}$ for $\rho < 0.25$ and $\vartheta_{\max} = \pi$ thereafter, and the sampled loop speed is $v_{\text{loop}} = 2.0(1-0.7\alpha) \pm 0.3$\,m/s. This branch exposes the policy to mid-maneuver states that on-policy rollouts rarely reach.

\textbf{Domain randomization.}
Per-environment perturbations are applied once at construction to improve sim-to-real robustness. The randomized parameters and ranges are listed in Table~\ref{tab:domrand}.

\section{Lessons Learned and Key Design Decisions}

The dense progress reward $r_{\text{prog}}$ was the most important design choice. Without it, the policy received almost no useful signal before the first gate pass. The $\tanh$-based shaping provided a stable long-range gradient without overwhelming the sparse gate bonus. Curriculum learning was the second key factor. Staged resets and mid-loop samples exposed the policy to states that on-policy rollouts rarely reached early in training.

Another useful finding was that the final policy often bypassed the intended vertical loop and flew around the powerloop gates instead. This detour was faster on our track. Disabling the crash penalty and action smoothness penalty also improved performance. The policy became more aggressive near gate edges, while the quadrotor dynamics already provided sufficient smoothness. The full curriculum and hyperparameters are listed in the Appendix.

A natural next step is to explore additional physics-informed reward heuristics. Examples include velocity alignment with the current gate direction, rewards that favor dynamically efficient high-speed flight, and terms that better capture the coupling among attitude, thrust, and future gate geometry. Such heuristics may encourage more time-efficient trajectories and further reduce lap time.

\section{Conclusion}

We presented a PPO policy for seven-gate drone racing in Isaac Sim. The final system combines dense global progress, sparse gate bonuses, a time penalty, a compact 36-dimensional observation, a staged reset curriculum, and domain randomization. In evaluation, the best policy completes three laps in 18.26\,s.

\bibliographystyle{IEEEtran}
\begin{thebibliography}{6}

\bibitem{schulman2017ppo}
J.~Schulman, F.~Wolski, P.~Dhariwal, A.~Radford, and O.~Klimov, ``Proximal policy optimization algorithms,'' \textit{arXiv preprint arXiv:1707.06347}, 2017.

\bibitem{song2021autonomous}
Y.~Song, M.~Steinweg, E.~Kaufmann, and D.~Scaramuzza, ``Autonomous drone racing with deep reinforcement learning,'' in \textit{Proc. IEEE/RSJ Int. Conf. Intell. Robots Syst. (IROS)}, 2021, pp.~1205--1212.

\bibitem{kaufmann2023champion}
E.~Kaufmann, L.~Bauersfeld, A.~Loquercio, M.~M{\"u}ller, V.~Koltun, and D.~Scaramuzza, ``Champion-level drone racing using deep reinforcement learning,'' \textit{Nature}, vol.~620, no.~7976, pp.~982--987, 2023.

\bibitem{song2023reaching}
Y.~Song, A.~Romero, M.~M{\"u}ller, V.~Koltun, and D.~Scaramuzza, ``Reaching the limit in autonomous racing: Optimal control versus reinforcement learning,'' \textit{Science Robotics}, vol.~8, no.~82, p.~eadg1462, 2023.

\bibitem{han2025reactive}
Z.~Han, X.~Huang, Z.~Xu, J.~Zhang, Y.~Wu, M.~Wang, T.~Wu, and F.~Gao, ``Reactive aerobatic flight via reinforcement learning,'' \textit{IEEE Robot. Autom. Lett.}, 2025.

\bibitem{ma2024eureka}
Y.~J.~Ma, W.~Liang, G.~Wang, D.-A.~Huang, O.~Bastani, D.~Jayaraman, Y.~Zhu, L.~Fan, and A.~Anandkumar, ``Eureka: Human-level reward design via coding large language models,'' in \textit{Proc. Int. Conf. Learn. Representations (ICLR)}, 2024.

\end{thebibliography}

\newpage
\appendix
\section{Curriculum Schedule}

\begin{table}[h]
\centering
\caption{Curriculum stages. $U$: number of unlocked reset targets, corresponding to gates $\{0,\dots,U-1\}$; $p_0$: probability of resetting at gate 0.}
\label{tab:curriculum}
\begin{tabular}{@{}ccccl@{}}
\toprule
\textbf{Stage} & $\bm{\rho}$ & $\bm{U}$ & $\bm{p_0}$ & \textbf{Focus} \\
\midrule
1 & $[0,\, 0.05)$ & 1 & 1.00 & Stable starts and gate 0 acquisition \\
2 & $[0.05,\, 0.12)$ & 2 & 0.55 & Introduce gate 1 resets \\
3 & $[0.12,\, 0.25)$ & 4 & 0.45 & Add powerloop entry and exit practice \\
4 & $[0.25,\, 0.50)$ & 6 & 0.40 & Extend to later gates and chicane entry \\
5 & $[0.50,\, 1.00]$ & 7 & 0.35 & Full-track training with residual gate 0 starts \\
\bottomrule
\end{tabular}
\end{table}

\section{Domain Randomization}

\begin{table}[h]
\centering
\caption{Domain randomization parameters. Nominal values are from the Crazyflie 2.x configuration.}
\label{tab:domrand}
\begin{tabular}{@{}lll@{}}
\toprule
\textbf{Parameter} & \textbf{Nominal} & \textbf{Range} \\
\midrule
Thrust-to-weight $T\!/\!W$ & 3.15 & $\times\,(1 \pm 0.05)$ \\
Aero drag $C_{d,xy}$ & $9.18\!\times\!10^{-7}$ & $\times\,\mathcal{U}(0.5,\,2.0)$ \\
Aero drag $C_{d,z}$ & $1.03\!\times\!10^{-6}$ & $\times\,\mathcal{U}(0.5,\,2.0)$ \\
PID $K_p$ (roll/pitch) & 250.0 & $\times\,(1 \pm 0.15)$ \\
PID $K_i$ (roll/pitch) & 500.0 & $\times\,(1 \pm 0.15)$ \\
PID $K_d$ (roll/pitch) & 2.5 & $\times\,(1 \pm 0.30)$ \\
PID $K_p$ (yaw) & 120.0 & $\times\,(1 \pm 0.15)$ \\
PID $K_i$ (yaw) & 16.7 & $\times\,(1 \pm 0.15)$ \\
Motor time constant $\tau_m$ & 0.005\,s & fixed \\
\bottomrule
\end{tabular}
\end{table}

\section{Training Hyperparameters}

\begin{table}[h]
\centering
\caption{PPO and network hyperparameters.}
\label{tab:ppo}
\begin{tabular}{@{}ll@{}}
\toprule
\textbf{Hyperparameter} & \textbf{Value} \\
\midrule
Discount factor $\gamma$ & 0.99 \\
GAE $\lambda$ & 0.95 \\
Clip ratio $\epsilon$ & 0.2 \\
Learning rate $\eta$ & $3\!\times\!10^{-4}$ (adaptive) \\
Target KL divergence $D_{\text{KL}}$ & 0.01 \\
Entropy coefficient & 0 \\
Max gradient norm & 1.0 \\
Steps per env per iteration & 24 \\
Number of parallel envs & 4096 \\
Mini-batches / Epochs & 4 / 5 \\
Batch size & $4096 \times 24 = 98{,}304$ \\
Total iterations $K$ & 1500 (reported run) \\
\midrule
Actor hidden layers & $[128, 128]$, ELU \\
Critic hidden layers & $[512, 256, 128, 128]$, ELU \\
Initial noise std & 1.0 \\
\midrule
Simulation rate & 500\,Hz ($dt=0.002$\,s) \\
Policy rate & 50\,Hz (decimation = 10) \\
Episode length & 30\,s ($\sim$1500 steps) \\
Action space & 4 (collective thrust + body-rate commands) \\
\bottomrule
\end{tabular}
\end{table}

\end{document}
